
\chapter{Discussion}\label{ch:discussion}

This chapter interprets the experimental results of Chapter~\ref{ch:experiments} in light of the design goals and the competitive landscape from Chapter~\ref{ch:litreview}. Section~\ref{sec:disc-security} discusses security effectiveness; Section~\ref{sec:disc-scanner} discusses scanner limitations; Section~\ref{sec:disc-positioning} discusses analytical positioning; Section~\ref{sec:disc-validity} connects the remaining validity bounds to future work.

\section{Security effectiveness after application compromise}\label{sec:disc-security}

The laboratory run supports the thesis claim that kernel-side confinement can be automated at the runtime layer without patching application source code. The Flask diagnostics endpoints remain vulnerable in the image, but scan-then-enforce prevents their exploitation at runtime. Under stock \texttt{runc}, the tested injection payloads effectively obtained an unconstrained root context inside the container namespace. Under hardened enforcement, the same payloads were blocked: the enforced seccomp profile omits process-creation and high-risk syscalls that the scan never observed, so the injected \texttt{sh -c} chain cannot run and follow-on reconnaissance (\texttt{id}, \texttt{hostname}, \texttt{uname}), sensitive reads, reverse shells, and mount probes fail as well.

This outcome must not be read as eliminating every remote-code-execution path in general. Other architectures, privileged or misconfigured deployments (\texttt{docker.sock} mounts, \texttt{--privileged}, excessive capabilities left in \texttt{config.json}), or applications whose harmful behaviour stays inside already-learned syscalls may still be exploitable even when profiles are present. On the laboratory stack, however, command injection and the tested follow-on payloads were prevented. Manually re-enabling \texttt{execve} in the seccomp profile restores a shell but still leaves most follow-up actions blocked, so the remaining layers add depth beyond the exec denial. This is a workload-specific result, not a guarantee for every attack scenario.

This separation matters for practitioners. Runtime-generated profiles are not a substitute for secure coding or web-application firewalls. They are a second line of defence that limits blast radius when the first line fails---the same role cluster-side recorders play when profiles are pinned to pods, but reachable from a developer workstation with one CLI flag.

The empirical comparison in this thesis is default \texttt{runc} vs.\ scan-then-enforce on a shared-kernel workload. Sandboxed runtimes remain relevant as an analytical alternative: they trade integration cost and overhead for hardware- or user-space-kernel isolation, whereas hardened \texttt{runc} stacks seccomp, AppArmor, and capabilities without leaving the OCI ecosystem.

\section{Scanner gaps and profile quality}\label{sec:disc-scanner}

The evaluation exposed three limitations that a reader should weigh when judging RQ1.

\textbf{AppArmor audit parsing.} Complain-mode learning is only useful if the collector sees every would-be denial. The implementation treats \texttt{ALLOWED} audit records as first-class evidence alongside \texttt{DENIED}, which is why writable calendar-storage paths are present in the generated profile and remain available after enforcement.

\textbf{File capabilities (limitation).} Behavioural capability tracing records kernel checks, not privileges embedded in executable files. \texttt{ping} is a common case: the raw-network capability on the binary does not necessarily produce \texttt{cap\_capable} events during a normal ping. Until file capabilities are merged from \texttt{security.capability} xattrs at exec time, operators need an explicit allow-list step---documented in the laboratory's \texttt{apply.sh}.

\textbf{Coverage equals exercised paths (limitation).} The generated profiles are only as complete as the traffic driven during the scan. The laboratory script covers the main UI and diagnostic flows but does not claim to hit cron jobs, certificate rotation, or rare error handlers. This limitation is shared with SPO, Inspektor Gadget, and every other dynamic profiler; the mitigation is to run the scan under the same integration tests that gate releases.

\section{Analytical positioning}\label{sec:disc-positioning}

Manual tuning and cluster recorders can, with sufficient effort, produce profiles at least as tight as the runtime scanner. The contribution of this work is lowering the operational floor: one flag, no cluster, no CRDs. Sandboxes address a different threat model (kernel isolation) and are not interchangeable with profile-based hardening. The empirical campaign in Chapter~\ref{ch:experiments} (Table~\ref{tab:perf-ch5}) shows gVisor 31--73\% below Docker on the same workloads while enforcing generated profiles on the hardened binary adds at most $\approx$5.5\% relative to the unprofiled bundle on Redis and less than 1\% on sysbench and iperf3.

\section{Validity bounds and future-work link}\label{sec:disc-validity}

The experiment is intentionally bounded: it uses one primary Flask workload for exploitation tests, qualitative succeed/block outcomes for RQ2, documented post-scan adjustments for file capabilities and writable paths, no privileged-container or \texttt{docker.sock} escape cases yet, and a VirtualBox guest without KVM for the first performance campaign. These bounds do not invalidate the observed result, but they are also not final claims about all container workloads.

For that reason, the remaining validity gaps are carried into Chapter~\ref{ch:conclusion} as future work: broader vulnerable workloads and escape-class tests (CVE-2019-5736, misconfigurations, post-RCE lateral movement), quantitative containment metrics beyond the Flask laboratory, automatic file-capability merge, release-to-release profile diffs, and a KVM-backed repeat of the performance campaign with full AppArmor enforcement on every workload.
